Quando discutimos entrada e saída de dados em C, frequentemente nos referimos 
a essa interação por meio de ``arquivos''. Embora esses ``arquivos'' não sejam, 
de fato, arquivos dentro do seu computador, essa é uma abstração 
útil para compreender o processo.

\label{o-que-sao-arquivos}Na prática, um "arquivo" é simplesmente um espaço na memória usado pelo 
sistema para armazenar informações de forma organizada. Uma analogia útil é 
imaginar que cada vez que o computador recebe uma informação, ela é escrita em 
um arquivo para armazenamento, e cada vez que uma informação precisa ser 
exibida, ela é escrita em um arquivo padrão que o terminal utiliza para 
determinar o que deve ser exibido na tela.

\subsubsection{Saída}

Além disso, é importante mencionar que os nomes oficiais, por assim a se dizer, 
utilizados para se referir a esses arquivos são: \texttt{stdin (entrada), 
stdout (saída)}.

Porém, um detalhe que vale ser mencionado, e que poucas pessoas se perguntam 
sobre é, existe algum outro tipo de arquivo que não sejam os de entrada e 
saída? E a resposta é que sim, esse arquivo seria o \texttt{stderr}. Que é 
unicamente designado para exibição de mensagens de erro e é completamente 
independente do \texttt{stdout}, permitindo assim, que haja uma distinção 
entre todo o conteúdo padrão que está sendo exibido no terminal e mensagens 
de erro.

Para demonstrar a utilização desses dois tipos de saída de dados diferentes 
vamos usar o seguinte exemplo:

\newpage
\begin{excode}
    #include <stdio.h>

    void print_error(void)
    {
        // Escreve a mensagem no arquivo `stdout'
        fprintf(stdout, "Essa é uma mensagem normal.");
    }

    void print_message(void)
    {
        // Escreve a mensagem no arquivo `stderr'
        fprintf(stderr, "Essa é uma mensagem de erro.");
    }

    int main(void)
    {
        print_message();
        print_error();
        return 0;
    }
\end{excode}

O código acima utiliza da função \texttt{fprintf} para direcionar o output de 
cada uma das mensagens de maneira mais específica. Diferente do \texttt{printf}, 
nessa outra função você deve escolher entre o \texttt{stdout e o stderr} antes 
de escolher qual mensagem você quer que seja exibida.

Ao compilar e executar o código você deve conseguir ver as duas mensagens 
sendo exibidas na tela da mesma forma, porém, as duas estão sendo exibidas de 
locais diferentes. Em sistemas Linux e Windows é possível filtrar o fluxo 
dessas mensagens da seguinte forma:

\begin{textbox}
    # Envia o conteúdo de stdout e stderr para 'output.txt'
    ./main > output.txt
    # Envia somente as mensagens de stderr para 'error.txt'
    ./main 2> error.txt
    # Envia somente as mensagens de stderr para 'error.txt'
    ./main > output.txt 2> error.txt
\end{textbox}

\subsubsection{Entrada}

Em C, a entrada de dados geralmente é realizada através do arquivo padrão de 
entrada, conhecido como \texttt{stdin}. Este arquivo é associado ao 
teclado, permitindo que o programa receba dados inseridos pelo usuário durante 
a execução.

Para ler dados do \texttt{stdin}, podemos usar funções de entrada como 
\texttt{scanf} e \texttt{fgets}. A função \texttt{scanf} é usada para ler 
entrada formatada, enquanto \texttt{fgets} é usada para ler uma linha de 
entrada como um texto.

Aqui está um exemplo de uso do \texttt{scanf} para ler um número inteiro:

\begin{excode}
    #include <stdio.h>

    int main(void)
    {
        int num;
        printf("Digite um número inteiro: ");
        scanf("\%d", &num);
        printf("Você digitou: \%d\n", num);
        return 0;
    }
\end{excode}

No exemplo acima, o programa solicita ao usuário que digite um número inteiro. 
O valor inserido pelo usuário é lido pelo \texttt{scanf} e armazenado na variável 
\texttt{num}, o qual é impressa na tela.

Se o programa precisar ler uma linha inteira de texto, pode-se usar a função 
\texttt{fgets}, como mostrado no exemplo a seguir:

\begin{excode}
    #include <stdio.h>

    int main(void)
    {
        char linha[100];
        printf("Digite uma frase: ");
        fgets(linha, 100, stdin);
        printf("Você digitou: %s\n", linha);
        return 0;
    }
\end{excode}

Neste exemplo, o programa solicita ao usuário que digite uma frase. A função 
\texttt{fgets} lê a linha de entrada e armazena na variável \texttt{linha}, 
a qual é então impressa na tela.

Esses são exemplos simples de como usar o \texttt{stdin} para entrada de dados 
em um programa em C. É importante lembrar de validar e manipular os dados de 
entrada adequadamente para garantir o bom funcionamento do programa.

\subsubsection{Exercícios}

\textbf{É de extrema importância que você realize os exercícios aqui elaborados para fixar o conteúdo.}

\begin{enumerate}
    \item {Escreva um programa que leia dois números inteiros do usuário usando 
        \texttt{scanf} e exiba a soma deles no \texttt{stdout}. Em seguida, 
        modifique o programa para que, caso o usuário digite um número negativo, 
        uma mensagem de erro seja enviada para o \texttt{stderr} em vez de 
        interromper o programa.}
    \item {Crie um programa que contenha duas funções que escreve uma mensagem 
        no \texttt{stdout} usando \texttt{fprintf} e uma que escreva a mensagem 
        no \texttt{stderr}. No \texttt{main}, chame ambas as funções com mensagens 
        diferentes. Compile e execute o programa redirecionando apenas o 
        \texttt{stderr} para um arquivo chamado \texttt{erros.txt}. Após isso, 
        analise o que você observa no terminal.}
    \item {Faça um programa que leia uma linha de texto completa (incluindo 
        espaços) usando \texttt{fgets} e, em seguida, exiba o texto lido caractere 
        por caractere (um por linha) utilizando um loop.}
\end{enumerate}
